\originalchapter{Doppler效应与膨胀宇宙的观测习题}
本章包含8.286 PS1三题与PS2六题，解答为本卷独立补写。声波需要空气参考系，光波需要局部测量者的四速度；宇宙学红移需要发射、接收时的尺度因子。三个问题有相似的频率比，却依赖不同的物理结构。

\src{8.286 PS1.1，声源与观察者同时运动。}
\begin{exercise}
声源相对空气速度$v_s$，观察者相对空气以$v_o$向相反方向离开声源。求$1+z=\lambda_o/\lambda_s$；这里以同一介质声速$u$把频率换为等效波长。
\end{exercise}
\begin{solution}
明确两者互相远离，并取声线从源到接收者的方向为正。源速度投影为$-v_s$，连续两波峰发射间隔$T_s$，在空气中波峰间距$(u+v_s)T_s$。接收者以$v_o$沿波的方向逃离，波峰相对其速度为$u-v_o$，所以
\[
 T_o=T_s\frac{u+v_s}{u-v_o},\qquad
 1+z=\frac{u+v_s}{u-v_o}.
\]
需要$v_o<u$才能依次接收到这条直达声线。任一方不动可恢复单方运动公式；小速度时$z\simeq(v_s+v_o)/u$，而精确表达仍各自依赖空气速度，不能只用相对速度替换分母。
\end{solution}

\src{8.286 PS1.2，横向Doppler，(a)--(e)。}
\begin{exercise}
源以速率$v$绕静止观察者作半径$R$圆周运动。(a)观察者静止于空气，$v,u\ll c$的声波红移是多少？
(b)若为光且$v$可相对论，求红移；
(c)反过来，由圆心向绕圆的接收者发光，求红移。
(d)源改为直线运动，光在源最靠近观察者时发射；
(e)接收者直线运动，光在接收者最靠近静止源时收到。分别求(d)、(e)红移。
\end{exercise}
\begin{solution}
(a) 圆心到发射点距离始终$R$，传播延迟始终$R/u$；连续波峰的接收时间差等于发射时间差，没有经典横向频移，$z=0$。

光的频率在实验室中为$\nu$，速度$\boldsymbol v$测量者所测$\nu_{\rm local}=\gamma\nu(1-\boldsymbol n\cdot\boldsymbol v/c)$，其中$\boldsymbol n$是该段光线在实验室的传播方向。这可由光波四矢量与四速度的内积直接得到。对横向发射，源测量频率$\nu_s=\gamma\nu$，静止接收者测$\nu_o=\nu$。因此(b)$1+z=\gamma$。

(c)实验室光径仍沿圆心至接收事件的半径，接收速度与其垂直，故$\nu_o=\gamma\nu_s$，$1+z=\gamma^{-1}$，是蓝移。发射与接收的角色不能在有加速的圆周世界线上直接交换后仍使用同一红移。

(d)在发射事件处，光线方向与源的直线速度垂直，局部频率关系仍给$1+z=\gamma$；(e)在接收事件处光方向与接收者速度垂直，所以$1+z=\gamma^{-1}$。关键是指定\emph{哪个事件}最靠近，若只说“看到最靠近时的源”会把传播延迟和发射位置混在一起。小速度展开分别为$z\simeq\pm v^2/(2c^2)$。
\end{solution}

\src{8.286 PS1.3，高速旋转木马，(a)--(b)，加分题。}
\begin{exercise}
四车同向以速率$v$绕半径$R$等间距转动，每车向其他三车发波。
(a)对$v,u\ll c$声波，求一辆车收到前车、后车、对面车波的红移；
(b)对光波及任意$v<c$，求同三种红移。空气取圆心实验室静止，忽略介质被车带动。
\end{exercise}
\begin{solution}
不能用同一时刻的连线代替发射事件至接收事件的声线／光线，但可直接利用稳恒旋转的对称性：将一条传播过程整体平移$\Delta t$并旋转$\omega\Delta t$，得到完全相同的下一条过程。因此同一传播分支的飞行时间是常数，发射与接收的实验室时间间隔相等。

还可用瞬时投影核查。设声线方向$\boldsymbol n$沿实际弦，两个事件的位置分别$\boldsymbol r_s,\boldsymbol r_o$，速度均为$\boldsymbol\omega\times\boldsymbol r$。于是
\[
 \boldsymbol n\cdot(\boldsymbol v_o-\boldsymbol v_s)
 =\boldsymbol n\cdot[\boldsymbol\omega\times(\boldsymbol r_o-\boldsymbol r_s)]=0 .
\]
经典频率比为$(u-\boldsymbol n\cdot\boldsymbol v_o)/(u-\boldsymbol n\cdot\boldsymbol v_s)=1$；相对论频率比另有$\gamma_o/\gamma_s$，两车速率相同也为1。故前、后、对面的每种情况都$z=0$。

这不表示彼此速度相同，而是端点沿传播方向的投影相等、两车钟的时间膨胀相等。声速条件应保证所讨论传播分支确实可用；若超声速导致多条到达分支或波峰聚集，需讨论信号形态，但每个稳恒分支的时间对称性没有改变。
\end{solution}

以下六题共同假设平坦尘埃宇宙$a(t)=bt^{2/3}$，今天$t_0$、发射$t_e$，无本动速度和吸收。原题历史例为UDFy--38135539的$z=8.55$；本章只按该2013题面参数计算，不把它称为今日最高红移纪录。尺度因子常数$b$应从所有可观测量消失。

\src{8.286 PS2.1，现今物理距离。}
\begin{exercise}
给$t_0,t_e$，求现在与所观测星系的物理距离$R_p$。
\end{exercise}
\begin{solution}
径向零测地线给共动距离$\chi=\int_{t_e}^{t_0}c\dd t/a(t)$，今天物理距离为
\[
 R_p=a_0\chi
 =3ct_0^{2/3}(t_0^{1/3}-t_e^{1/3})
 =3ct_0[1-(t_e/t_0)^{1/3}].
\]
它是在同一$t_0$切片摆放共动尺的距离，不是$c(t_0-t_e)$。光在飞行中经历膨胀，所以两种量不相等。
\end{solution}

\src{8.286 PS2.2，发射时刻。}
\begin{exercise}
用$t_0,t_e$表示$z$，求$z=8.55$时$t_e/t_0$。
\end{exercise}
\begin{solution}
共动发、收者测的光频率反比$a$，故
\[
 1+z=\frac{a_0}{a_e}=(t_0/t_e)^{2/3},\qquad
 \frac{t_e}{t_0}=(1+z)^{-3/2}=0.0338840.
\]
这个很小的比例属于尘埃模型；实际早期辐射、暗能量会改变红移与年龄的关系，不能照此当作真实星系年代测定。
\end{solution}

\src{8.286 PS2.3，以Hubble常数表示距离。}
\begin{exercise}
用$H_0,z$表示$R_p$，对$H_0=67\ \mathrm{km\,s^{-1}\,Mpc^{-1}},z=8.55$求Mpc及光年。
\end{exercise}
\begin{solution}
$H_0=2/(3t_0)$，代入前两题：
\[
 R_p=\frac{2c}{H_0}\left[1-(1+z)^{-1/2}\right].
\]
实际求值$6053.19$ Mpc，或$1.97429\times10^{10}$光年。低红移展开$R_p\simeq cz/H_0$恢复Hubble近似；高红移极限$2c/H_0=3ct_0$是本模型今天的粒子视界。
\end{solution}

\src{8.286 PS2.4，退行速率。}
\begin{exercise}
求星系今天物理距离的变化率$\dot R_p$，用$c,z$表示，并求$z=8.55$时与$c$的比。
\end{exercise}
\begin{solution}
星系共动坐标固定，$\dot R_p=H_0R_p$，所以
\[
 \frac{\dot R_p}{c}=2[1-(1+z)^{-1/2}]=1.352815.
\]
原题把超光速退行称为“违反SR”，准确说法是：这不是一个局部惯性系中两个邻近物体的相对速度，因而不受SR局部速度上限直接约束。沿光路每个局部共动观察者仍测光速$c$，GR并没有要求局部物体超光速。
\end{solution}

\src{8.286 PS2.5，角大小，含极值比较。}
\begin{exercise}
发射时球状星系的物理直径为$w$，$\theta\ll1$。用$w,z,H_0,c$求今天角直径，与静止欧氏空间中相同直径、距离为今天$R_p$的角直径比较，求角直径最小的红移。
\end{exercise}
\begin{solution}
小物体两边的共动横向间距为$w/a_e$，所以$\theta=(w/a_e)/\chi=w/D_A$，其中
\[
 D_A=\frac{R_p}{1+z}
 =\frac{2c}{H_0}\frac{1-(1+z)^{-1/2}}{1+z}.
\]
因此$\theta=(1+z)w/R_p$，比题目欧氏对照角度大$1+z$倍。令$y=\sqrt{1+z}$，$D_A=(2c/H_0)(y-1)/y^3$，导数正比$3-2y$，极大值在$y=3/2$，即$z=1.25$；这里$\theta$最小。更高红移看到的发射切片更早，角径距反而减小。上述固定$w$比较不包括真实星系自身尺寸的演化。
\end{solution}

\src{8.286 PS2.6，辐射通量。}
\begin{exercise}
星系发射时总功率$P$，各向同性地以光子辐射，求今天地球收到的垂直能流$J$。
\end{exercise}
\begin{solution}
将星系设为共动原点，光子今天分布在面积$4\pi R_p^2$球面。每光子能量减小$1+z$倍，两批光子的到达时间间隔比发射间隔延长$1+z$倍，因此
\[
 J=\frac{P}{4\pi R_p^2(1+z)^2}
 =\frac{PH_0^2}{16\pi c^2(1+z)^2[1-(1+z)^{-1/2}]^2}.
\]
识别$D_L=(1+z)R_p=(1+z)^2D_A$。两次红移因子分别对应能量和到达率，不能只保留一个；也不能把正在传播的光球面积错用为发射时的物理面积。
\end{solution}
